\section{Internalization audit and the Lean formalization path}
\label{sec:formalization}

\subsection{Why the proof fits quantifier-free induction}

It is insufficient to say that an algorithm is computable and therefore
available in iPRA. The algorithms used here have specified primitive
recursive bounds, and their correctness proofs can be organized as
quantifier-free finite invariants. The following table records what is
actually needed.

\begin{longtable}{@{}p{0.25\linewidth}p{0.32\linewidth}p{0.34\linewidth}@{}}
\toprule
Construction & Numerical data and bound & Induction invariant\\
\midrule
\endhead
Expression syntax & Tagged node array, child indices below the current index & Every processed node has the required tag and well-formed children.\\
Taylor approximations & Explicit finite cutoff in input magnitude and precision & Partial-sum recurrence and factorial/geometric tail inequalities.\\
Expression evaluation & Global magnitude and precision bounds from Appendix~\ref{app:primitives} & Bounded node table of rational values and error budgets.\\
Coarse rounding & One rational approximation with error $1/8$; rational floor & Quotient/remainder inequalities and $5/8$ distance bound.\\
Polynomial majorant & Finite coefficient and exponent lists & Telescoping product identity and accumulated rational error bound.\\
Density, phase search & Indices at most $K,N_0,N_1$ in Section~\ref{sec:density} & Finite grid comparisons, signed rational endpoint values, and adjacent-sample bound.\\
Density, dimension loop & Exactly $n$ nested rational intervals & Certificate table of centers, radii, Lipschitz bounds, and interval inclusions.\\
Square targets & Fixed number of bisections determined by $\rho$ and $M$ & Rational endpoints satisfy $l^2\le m_i\le u^2$ and the interval-length recurrence.\\
Source-proof translation & One fixed finite EFA proof & Each translated line is derived from translated axioms or earlier lines.\\
Machine simulation & Length of a supplied finite run certificate & Local configurations agree under the wrapper simulation.\\
\bottomrule
\end{longtable}

For example, the density construction need not induct on a statement
quantifying over all real $t$ in an interval. At stage $i$ store a rational
center $c_i$, a rational radius $r_i$, the integer Lipschitz bound $D_i$,
and a rational approximation to $h_i(c_i)-a_i$. A finite inequality
certifies that its absolute value, its evaluation error, and $D_ir_i$
sum to less than the tolerance. Subsequent intervals satisfy rational
containment tests. A primitive recursive Boolean verifier checks the
whole finite table. Induction is applied to that verifier's equation,
with the input tuple and requested precision as parameters.

The real notation in the analytic sections is a readable organization of
these rational error arguments. To translate an exact equality of
continuous operations, introduce an arbitrary final precision, choose the
explicitly sufficient auxiliary precisions, and prove the corresponding
rational inequality. The quantifier over the final precision is introduced
after the finite argument. This does not induct on an unbounded formula
asserting truth of a real identity.

The statements used by the final reduction can be kept especially simple.
Let $\operatorname{Root}(p,w)$ be the primitive recursive check that $w$
is a natural root tuple of polynomial code $p$. Let
$\operatorname{Obs}(a,z)$ decode a rational input and precision and perform
the strict test \eqref{eq:nz}. For primitive recursive witness maps
$\operatorname{forth}$ and $\operatorname{back}$, prove the universal
closures of
\begin{align*}
 \operatorname{Root}(p,w)&\to
   \operatorname{Obs}(E(p),\operatorname{forth}(p,w)),\\
 \operatorname{Obs}(E(p),z)&\to
   \operatorname{Root}(p,\operatorname{back}(p,z)).
\end{align*}
The matrices contain only decidable numerical relations. Existential
introduction/elimination and intuitionistic propositional logic then
produce the equivalence of witness existence and its negative version.

\subsection{What ProveIt supplies}
\label{sec:repo}

The repository was inspected through the GitHub connector at the pinned
commit
\begin{center}\small\code{eb8eb6653c85aa44f1e18dd4efa243f41b5a9e98}.\end{center}
The principal relevant interfaces are these \cite{proveit}.

\begin{enumerate}
\item \path{ConciseIPRA/Lean/ConciseIPRA/PR.lean}\par
This file, under \path{Logic/PrimitiveRecursiveArithmetic/}, defines the actual object
language: \code{PR k}, \code{PTm n}, \code{PRAAx}, and the quantifier-free
induction helper \code{PRA.qfInd}. The namespace is
\code{LeanProofs.ConciseIPRA}.
\item \code{Derivations/Lean/IPRADerivations/Core/Search.lean} supplies
object-level bounded search, including \code{search\_spec},
\code{search\_minimal}, and \code{search\_found}. It returns the bound on
failure. This is directly suitable for the finite grids in the density
proof, once their success bounds have been established.
\item The derivations library supplies arithmetic, order, pairing, finite
sums, and rational arithmetic. These are useful bases for the certificate
calculations, rather than reasons to introduce new axioms.
\item \path{Computability/HilbertTenthProblem/Lean/MRDP.md} and
\path{Diophantine/Common/MRDPCore.lean} under that directory supply an
ordinary Lean/Mathlib MRDP theorem. Its proof uses Diophantine
exponentiation, CRT encoding, and a bounded computation evaluator.
These are mathematical proof templates, not a preexisting object-level
iPRA representation theorem.
\end{enumerate}

The repository's own \code{UNIFICATION-PLAN.md} explicitly records that
boundary and treats general internal MRDP as separate work. Its README
also distinguishes checked \code{Prf PRAAx} derivations from delivered
research packets that have not been verified. Accordingly, this report
does not use one of those packets as a checked foundational theorem.
The full audit, precise source paths, and provenance are in
\code{evidence/repository-audit.txt}.

\subsection{Supplied Lean source and its exact status}

The archive's \code{lean/Syntax.lean} defines the elementary-expression
AST, substitution for composition, subtraction by multiplication by
$-1$, natural powers, and the operation $A\mapsto |A|-A$.
The file \code{lean/LogicalCore.lean} contains fully written constructive
proof terms for:
\begin{itemize}
\item negative calibration from a double-negated equivalence;
\item existence and absence transport from finite witness maps;
\item the two branches of the diagonal contradiction;
\item nonexistence of a total correct decider on the compiled domain.
\end{itemize}
Both files import only \code{Init}. They contain no axiom declarations
or proof holes. They were inspected but not compiled: no Lean executable
was available in the execution environment. No kernel-checking claim is
made for these new files.

The main logical theorem has explicit structure parameters for run
semantics, determinism, compilation, simulation, and calibration.
Such parameters expose the arithmetic and analytic obligations; they do
not discharge them. In particular, \code{Framework.reduction} is a
hypothesis of the parametric logical theorem. It must eventually be
replaced by the derived object-language statement
\eqref{eq:calibration}.

The essential logical step is small enough to display in full:
\begin{lstlisting}
theorem negIffOfDoubleNegIff {H S : Prop}
    (h : (Not (Not H)) <-> (Not (Not S))) :
    (Not H) <-> (Not S) :=
  Iff.intro
    (fun notH s => (h.mpr (fun notS => notS s)) notH)
    (fun notS h0 => (h.mp (fun notH => notH h0)) notS)
\end{lstlisting}
This is an ambient Lean proposition illustrating the proof transformation.
It is not yet an object-language proof value.

\subsection{The final object-level endpoint}

After defining the coded predicates, create the closed formula
\begin{lstlisting}
richardsonSentence : Fm PTm 0
\end{lstlisting}
which is the explicit first-order sentence
$\forall d\,\neg\Dec(d)$. The intended completed endpoint is
\begin{lstlisting}
theorem richardson :
  ([] : List (Fm PTm 0)) |- [PRAAx] richardsonSentence
\end{lstlisting}
Here the displayed ASCII turnstile is explanatory notation for the
repository's actual Unicode derivation notation
\code{Prf PRAAx}; the accompanying plan states the exact syntax.
It is not a claim that the code block above is a ready-to-compile module.

A practical completion order is:
\begin{enumerate}
\item Implement PR codes for syntax, rational operations, finite series,
the numerical observation test, and machine traces. Reuse the existing
bounded-search and arithmetic derivations.
\item Derive rational sine identities and bounds, the evaluator error
theorems, and the coarse-rounding lemma.
\item Derive the polynomial majorant, integer gap, density certificate
constructor, and the two witness-map correctness formulas.
\item Formalize the fixed source arithmetic proof and its negative
translation, or provide the same negative representation directly in
\code{Prf PRAAx}. This is the principal deep remaining formalization task.
\item Derive the wrapper trace laws and combine the internal formulas
with the two-branch logical proof.
\item Compile the object theorem, inspect its assumptions and imports,
and verify that the theorem returns a \code{Prf PRAAx} value for the
intended closed formula.
\end{enumerate}

The detailed file-level obligations and numerical signatures are included
in \path{lean/FormalizationPlan.txt}. This is a concrete path to completion;
the report does not substitute a classical \code{Real}-valued theorem or
an assumed reduction interface for the required internal endpoint.
